% !TeX root = ../sections/03_solutions.tex

\begin{solution}
	大問2では，中国剰余定理
	\[
	\mathbb{Z}/15\mathbb{Z}
	\cong
	\mathbb{Z}/3\mathbb{Z}
	\oplus
	\mathbb{Z}/5\mathbb{Z}
	\]
	を具体的に用いる。
	
	ここで，
	\[
	\mathbb{Z}/3\mathbb{Z}
	\oplus
	\mathbb{Z}/5\mathbb{Z}
	\]
	の元
	\[
	(r,r)
	\]
	とは，正確には
	\[
	(r \bmod 3,\ r \bmod 5)
	\]
	を表していると考える。
	したがって，
	\[
	(r,r)=(s,t)
	\]
	とは，
	\[
	r\equiv s \pmod{3},
	\qquad
	r\equiv t \pmod{5}
	\]
	が同時に成り立つという意味である。
	
	\begin{enumerate}
		\item[(0)]
		\[
		1=3a+5b
		\]
		を満たす整数 \(a,b\) を一組求める。
		
		まず，
		\[
		1=6-5
		\]
		である。ここで，
		\[
		6=3\cdot 2,
		\qquad
		5=5\cdot 1
		\]
		なので，
		\[
		1=3\cdot 2-5\cdot 1
		\]
		と書ける。
		
		したがって，
		\[
		1=3\cdot 2+5\cdot(-1)
		\]
		である。
		
		よって，
		\[
		\boxed{
			a=2,\quad b=-1
		}
		\]
		は条件を満たす。
		
		実際，
		\[
		3a+5b
		=
		3\cdot 2+5\cdot(-1)
		=
		6-5
		=
		1
		\]
		である。
		
		\item[(1)]
		\[
		\mathbb{Z}/3\mathbb{Z}
		\oplus
		\mathbb{Z}/5\mathbb{Z}
		\]
		において，
		\[
		(5b,5b)=(1,0)
		\]
		であることを示す。
		
		(0) で
		\[
		b=-1
		\]
		と取ったので，
		\[
		5b=5\cdot(-1)=-5
		\]
		である。
		
		したがって，
		\[
		(5b,5b)=(-5,-5)
		\]
		である。
		
		ここで，第一成分は \(\mathbb{Z}/3\mathbb{Z}\) で考える。
		\[
		-5\equiv 1 \pmod{3}
		\]
		である。実際，
		\[
		-5-1=-6
		\]
		であり，
		\[
		-6=3\cdot(-2)
		\]
		だから，
		\[
		3\mid (-5-1)
		\]
		である。
		よって，
		\[
		-5\equiv 1\pmod{3}
		\]
		である。
		
		次に，第二成分は \(\mathbb{Z}/5\mathbb{Z}\) で考える。
		\[
		-5\equiv 0 \pmod{5}
		\]
		である。実際，
		\[
		-5-0=-5
		\]
		であり，
		\[
		5\mid -5
		\]
		だからである。
		
		したがって，
		\[
		(-5,-5)=(1,0)
		\]
		である。
		
		ゆえに，
		\[
		\boxed{
			(5b,5b)=(1,0)
		}
		\]
		が成り立つ。
		
		\item[(2)]
		\[
		\mathbb{Z}/3\mathbb{Z}
		\oplus
		\mathbb{Z}/5\mathbb{Z}
		\]
		において，
		\[
		(3a,3a)=(0,1)
		\]
		であることを示す。
		
		(0) で
		\[
		a=2
		\]
		と取ったので，
		\[
		3a=3\cdot 2=6
		\]
		である。
		
		したがって，
		\[
		(3a,3a)=(6,6)
		\]
		である。
		
		ここで，第一成分は \(\mathbb{Z}/3\mathbb{Z}\) で考える。
		\[
		6\equiv 0 \pmod{3}
		\]
		である。実際，
		\[
		6-0=6
		\]
		であり，
		\[
		3\mid 6
		\]
		だからである。
		
		次に，第二成分は \(\mathbb{Z}/5\mathbb{Z}\) で考える。
		\[
		6\equiv 1 \pmod{5}
		\]
		である。実際，
		\[
		6-1=5
		\]
		であり，
		\[
		5\mid 5
		\]
		だからである。
		
		したがって，
		\[
		(6,6)=(0,1)
		\]
		である。
		
		ゆえに，
		\[
		\boxed{
			(3a,3a)=(0,1)
		}
		\]
		が成り立つ。
		
		\item[(3)]
		\[
		\mathbb{Z}/3\mathbb{Z}
		\oplus
		\mathbb{Z}/5\mathbb{Z}
		\]
		において，
		\[
		(4\times 3a+2\times 5b,\ 4\times 3a+2\times 5b)
		=
		(2,4)
		\]
		であることを示す。
		
		(1) より，
		\[
		(5b,5b)=(1,0)
		\]
		である。また，(2) より，
		\[
		(3a,3a)=(0,1)
		\]
		である。
		
		したがって，
		\[
		2(5b,5b)=2(1,0)=(2,0)
		\]
		であり，
		\[
		4(3a,3a)=4(0,1)=(0,4)
		\]
		である。
		
		ゆえに，
		\[
		4(3a,3a)+2(5b,5b)
		=
		(0,4)+(2,0)
		=
		(2,4)
		\]
		である。
		
		左辺は成分ごとの和であるから，
		\[
		4(3a,3a)+2(5b,5b)
		=
		(4\cdot 3a,\ 4\cdot 3a)
		+
		(2\cdot 5b,\ 2\cdot 5b)
		\]
		である。
		したがって，
		\[
		4(3a,3a)+2(5b,5b)
		=
		(4\cdot 3a+2\cdot 5b,\ 4\cdot 3a+2\cdot 5b)
		\]
		である。
		
		よって，
		\[
		\boxed{
			(4\times 3a+2\times 5b,\ 4\times 3a+2\times 5b)
			=
			(2,4)
		}
		\]
		が成り立つ。
		
		なお，実際に数値を入れて確認すると，
		\[
		a=2,\quad b=-1
		\]
		であるから，
		\[
		4\times 3a+2\times 5b
		=
		4\cdot 3\cdot 2+2\cdot 5\cdot(-1)
		\]
		である。よって，
		\[
		4\times 3a+2\times 5b
		=
		24-10
		=
		14
		\]
		である。
		
		このとき，
		\[
		14\equiv 2\pmod{3}
		\]
		であり，
		\[
		14\equiv 4\pmod{5}
		\]
		である。
		したがって，
		\[
		(14,14)=(2,4)
		\]
		であり，確かに主張が成り立つ。
		
		\item[(4)]
		\(3\) で割った余りが \(2\) となり，\(5\) で割った余りが \(4\) であるような整数を求める。
		
		求める整数を \(n\) とする。この条件は，
		\[
		n\equiv 2\pmod{3}
		\]
		かつ
		\[
		n\equiv 4\pmod{5}
		\]
		であることを意味する。
		
		これは，
		\[
		(n,n)=(2,4)
		\]
		が
		\[
		\mathbb{Z}/3\mathbb{Z}
		\oplus
		\mathbb{Z}/5\mathbb{Z}
		\]
		で成り立つことと同じである。
		
		(3) より，
		\[
		(4\times 3a+2\times 5b,\ 4\times 3a+2\times 5b)
		=
		(2,4)
		\]
		であった。したがって，
		\[
		n=4\times 3a+2\times 5b
		\]
		と取れば条件を満たす。
		
		ここで，
		\[
		a=2,\quad b=-1
		\]
		を代入すると，
		\[
		n
		=
		4\cdot 3\cdot 2
		+
		2\cdot 5\cdot(-1)
		\]
		である。
		よって，
		\[
		n=24-10=14
		\]
		である。
		
		実際，
		\[
		14=3\cdot 4+2
		\]
		なので，
		\[
		14\equiv 2\pmod{3}
		\]
		である。また，
		\[
		14=5\cdot 2+4
		\]
		なので，
		\[
		14\equiv 4\pmod{5}
		\]
		である。
		
		したがって，\(14\) は条件を満たす。
		
		さらに，条件を満たす整数全体を求める。
		もし \(n\) と \(14\) がどちらも条件を満たすならば，
		\[
		n\equiv 14\pmod{3}
		\]
		かつ
		\[
		n\equiv 14\pmod{5}
		\]
		である。
		
		したがって，
		\[
		3\mid(n-14)
		\]
		かつ
		\[
		5\mid(n-14)
		\]
		である。
		
		ここで，\(3\) と \(5\) は互いに素であるから，
		\[
		15\mid(n-14)
		\]
		である。
		
		よって，ある整数 \(t\in\mathbb{Z}\) が存在して，
		\[
		n-14=15t
		\]
		と書ける。したがって，
		\[
		n=14+15t
		\]
		である。
		
		逆に，
		\[
		n=14+15t
		\qquad
		(t\in\mathbb{Z})
		\]
		とすると，
		\[
		n\equiv 14\equiv 2\pmod{3}
		\]
		であり，
		\[
		n\equiv 14\equiv 4\pmod{5}
		\]
		である。
		したがって，すべての \(14+15t\) は条件を満たす。
		
		以上より，求める整数全体は
		\[
		\boxed{
			n=14+15t
			\qquad
			(t\in\mathbb{Z})
		}
		\]
		である。
	\end{enumerate}
\end{solution}